\documentclass[../../thesis_main.tex]{subfiles}
\graphicspath{{\subfix{../../figures/}}}

\begin{document}

    \chapter*{Abstract}
    This work presents the design, implementation and characterization of an Optical Dipole Trap (ODT) for capturing cold $^{87}$Rb atoms from a Magneto-Optical Trap (MOT) and optical molasses setup. 

    The ODT is realized by coupling 1064~nm laser light into a Hollow-Core Photonic Crystal Fiber (HCPCF). Due to the inherently multimodal nature of these fibers, light injection is highly sensitive to misalignment. Monitoring the near-field at the fiber output was found to be the most effective alignment procedure. For this purpose, a near-field imaging system was built to optimize the light injection, achieving a coupling efficiency of $96\%$ with an $LP$-like fundamental mode profile in test fibers. Subsequent characterization of an in-vacuum fiber revealed mode degradation and reduced transmission efficiency, hypothesized to result from mechanical stress induced by the vacuum sealing mechanisms.

    The final temperature of the optical molasses was minimized by systematically tuning the cooling beam detuning, the Acousto-Optic Modulator (AOM) ramp durations and amplitudes controlling the beam intensity, and the compensation coil currents used to cancel external magnetic fields. A final $xy$-averaged minimum temperature of $T = (9.6 \pm 0.1)~\mu\text{K}$ was achieved.

    Finally, the loading efficiency from the molasses into the ODT was optimized. After the MOT loading phase, the compensation coils were used to shift the atomic cloud's position, increasing the spatial overlap with the ODT beam. A loading efficiency of $(0.160 \pm 0.005)\%$ was achieved. This result is of the same order of magnitude as the expected value of $(0.41\pm0.04)\%$, calculated using the non-shallow trap model proposed in this thesis. The obtained efficiency is sufficient for the subsequent realization of an optical conveyor belt to effectively load the atoms into the fiber's core.

    This work establishes a reliable experimental starting point for introducing the second, counter-propagating 1064~nm beam, necessary to realize the optical conveyor belt.

\end{document}